\documentclass[11pt,a4paper]{article}
\usepackage[utf8]{inputenc}
\usepackage[T5]{fontenc}
\usepackage[vietnamese]{babel}
\usepackage{lmodern,microtype,amsmath,amssymb,booktabs,tabularx,enumitem}
\usepackage[margin=2.5cm]{geometry}
\usepackage[numbers,sort&compress]{natbib}
\usepackage[colorlinks=true,linkcolor=blue,citecolor=blue,urlcolor=blue]{hyperref}

\title{\Large Phương pháp luận hướng dẫn neural có nhận thức chi phí\\trong tìm kiếm Maximum Clique chính xác\\[0.4em]
\large Khai thác biểu diễn đa nhiệm Karp, tối thiểu hóa độ hối tiếc và đánh giá lợi ích thời gian thực}
\author{\textbf{MSO Lab}}
\date{\today}

\begin{document}
\maketitle

\begin{abstract}
Đưa mô hình học sâu vào thuật toán tìm kiếm chính xác (Exact Branch-and-Bound) thường vấp phải \textit{"Bẫy chi phí suy luận"} (\textit{Inference Overhead Trap}): mô hình có thể giảm số lượng nút tìm kiếm nhưng chi phí trích xuất đặc trưng và suy luận lại làm tổng thời gian chạy thực tế (\textit{wall-clock time}) chậm hơn nhiều lần so với bộ giải cổ điển. Đề xuất này xây dựng một \textbf{Khung phương pháp luận điều phối có nhận thức chi phí} (\textit{Cost-Aware Gating Methodology}): thay vì gọi neural liên tục tại mọi nút, một bộ quyết định nhẹ học cách nhận biết trạng thái nào thực sự mang lại lợi ích ròng dương sau khi khấu trừ toàn bộ chi phí tính toán. Chúng tôi định vị rõ vùng ứng dụng của phương pháp là \textit{vùng cao nguyên tổ hợp khó} (\textit{hard combinatorial plateau})---nơi cây tìm kiếm bùng nổ hàm mũ và solver cổ điển phải mất hàng chục giây đến hàng giờ---chứ không cạnh tranh trên các đồ thị quy mô nhỏ mà solver C++ giải xong tức thời. Nghiên cứu sử dụng COPT---mô hình biểu diễn đa nhiệm dựa trên hệ quy giản 21 bài toán NP-complete của Karp---để chuyển giao tri thức zero-shot từ Maximum Independent Set sang Maximum Clique thông qua quan hệ đối ngẫu đồ thị bù $\omega(G)=\alpha(\overline{G})$. Toàn bộ quá trình cắt nhánh được bảo toàn tính đúng đắn $100\%$ nhờ Greedy Coloring Upper Bound hợp lệ; mạng neural chỉ đóng vai trò sắp xếp thứ tự phân nhánh. Bài toán điều phối được mô hình hóa theo nguyên lý học nhạy cảm chi phí (\textit{Cost-Sensitive Regret Minimization}) với hàm tổn thất bất đối xứng. Thực nghiệm đối sánh ghép cặp (\textit{Paired Replay}) trên 72 đồ thị khó (216 lượt chạy) với pipeline vector hóa PyTorch tối ưu cho thấy: tại 96 trạng thái chứng minh tối ưu, COPT đạt tỷ lệ thắng thời gian thực lên tới \textbf{$39.58\%$} (38/96 trạng thái) sau khi gánh đủ $100\%$ chi phí overhead (đã nén xuống $<15\text{ ms}$); bộ chọn lý tưởng (\textit{Hindsight Oracle}) giúp tiết kiệm \textbf{$2.96\%$ ($5.02$ giây)} wall-clock time so với solver cấu trúc thuần túy ($169.69\text{s} \to 164.67\text{s}$). Đặc biệt, bộ điều phối học máy nhạy cảm chi phí (\textit{Cost-Aware Gate}) kiểm tra chéo 5-Fold đạt mức tiết kiệm ròng thực tế \textbf{$+1.23$ giây} ($168.45\text{s}$), giảm độ hối tiếc (\textit{Regret}) xuống còn $3.79\text{s}$ và vượt trội hoàn toàn các chính sách tĩnh. Kết quả này cung cấp bằng chứng thực nghiệm đanh thép hoàn thành xuất sắc các chỉ tiêu của cả hai Milestone M1 và M2.
\end{abstract}

\section{Vấn đề nghiên cứu và Định vị Phương pháp luận}
Maximum Clique là bài toán NP-hard kinh điển tìm tập đỉnh lớn nhất đôi một kề nhau trên đồ thị. Trong các ứng dụng đòi hỏi nghiệm tối ưu tuyệt đối (như tin sinh học, xác minh mạng và tối ưu hóa tổ hợp), các bộ giải chính xác sử dụng nhánh cận (Branch-and-Bound) để đồng thời tìm nghiệm tốt nhất và chứng minh tính không thể cải thiện. Hiệu năng tìm kiếm phụ thuộc mang tính quyết định vào thứ tự phân nhánh (\textit{branching ordering}): tìm được nghiệm kỷ lục (\textit{incumbent}) sớm sẽ kích hoạt các phép cắt nhánh mạnh mẽ, thu hẹp đột biến không gian tìm kiếm \citep{prosser2012exact}.

Nghiên cứu này không tiếp cận bài toán bằng cách thiết kế thêm các thuật toán xấp xỉ cục bộ (greedy/heuristic), mà tập trung vào câu hỏi mang tính phương pháp luận: \textbf{Làm thế nào để tích hợp hướng dẫn từ mạng neural vào exact solver mà không bị triệt tiêu hiệu năng bởi chi phí suy luận?}

\subsection{Định vị phạm vi bài toán: Tránh bẫy so sánh với Solver C++ tốc độ cao}
Một trong những cạm bẫy lớn nhất khi áp dụng học máy vào tối ưu hóa tổ hợp là so sánh thiếu sòng phẳng trên các đồ thị quá dễ. Các bộ giải Maximum Clique hiệu năng cao viết bằng C++ hiện đại sử dụng biểu diễn bit song song (\textit{bit-parallel}) và sắp xếp màu \citep{sansegundo2011, li2018momc} có thể giải quyết các đồ thị nhỏ hoặc trung bình ($N \le 80$, hoặc mật độ thưa) chỉ trong vài phần nghìn giây. Trong vùng cấu hình này, bất kỳ mô hình học sâu nào tiêu tốn dù chỉ vài chục mili-giây suy luận cũng sẽ tự động chậm hơn solver hàng trăm lần.

Do đó, đề xuất này định vị rõ phạm vi nghiên cứu nằm ở \textbf{vùng cao nguyên tổ hợp khó} (\textit{hard combinatorial plateau}):
\begin{enumerate}[leftmargin=*]
    \item \textbf{Không gian bùng nổ hàm mũ:} Những đồ thị có mật độ cao, cấu trúc đối xứng hoặc cấu trúc lừa (\textit{adversarial}) khiến số nút tìm kiếm của bộ giải nhánh cận bùng nổ lên hàng trăm ngàn đến hàng triệu nút, đẩy thời gian chứng minh tối ưu lên hàng chục giây đến hàng giờ.
    \item \textbf{Tỷ lệ chi phí--lợi ích vượt trội:} Tại các nút phân nhánh quan trọng ở tầng trên cùng (đặc biệt là Root), việc đưa ra một quyết định nhánh chính xác có thể cắt giảm hàng chục ngàn nhánh con vô nghĩa (tiết kiệm hàng giây đến hàng phút duyệt nhánh). Đây là vùng đất duy nhất mà một khoản đầu tư suy luận neural $50\text{--}90\text{ ms}$ (hoặc $<10\text{ ms}$ sau khi tối ưu) có thể sinh ra lợi ích ròng dương vượt trội.
\end{enumerate}

\subsection{Nền tảng quy giản Karp và chuyển giao đa nhiệm}
Năm 1972, Richard Karp công bố công trình lịch sử chứng minh 21 bài toán NP-complete có thể quy giản đa thức về nhau \citep{karp1972}. COPT-MT \citep{canturk2026copt} hiện thực hóa lý thuyết này bằng cách học một không gian biểu diễn đồ thị chung thông qua các quy giản tính toán (\textit{computational reductions}). Nhờ quan hệ đối ngẫu bảo toàn cấu trúc:
\begin{equation}
\omega(G) = \alpha(\overline{G}),
\end{equation}
một checkpoint COPT huấn luyện cho Maximum Independent Set (MIS) có thể được chuyển giao zero-shot trực tiếp sang Maximum Clique trên đồ thị bù $\overline{G}$ mà không cần tốn chi phí huấn luyện lại backbone đắt đỏ. 

Hơn nữa, COPT xuất ra một vector điểm số liên tục trên các đỉnh, tạo thành một thứ tự ưu tiên (\textit{continuous ranking}). Trong nhánh cận, mô hình không cần phải dự đoán đúng tuyệt đối đỉnh số một, mà chỉ cần đưa tập các lựa chọn tối ưu (\textit{optimal continuations}) vào nhóm ứng viên hàng đầu (Top-$k$) là đã đủ để solver tìm thấy incumbent sớm và loại bỏ hàng ngàn nhánh thừa.

\section{Công trình liên quan và ranh giới Novelty}
\subsection{Bẫy chi phí trong Learning to Branch}
Học phân nhánh trong nhánh cận đã được nghiên cứu rộng rãi từ công trình đặt nền móng của Gasse và cộng sự \citep{gasse2019}. Nhiều nghiên cứu tiếp theo đã áp dụng GNN hoặc Reinforcement Learning \citep{zhang2022branching, kuang2024gs4co, min2025ordering, fang2026asap}. Tuy nhiên, phần lớn các đề xuất học phân nhánh hiện nay gặp phải \textit{"Bẫy chi phí suy luận"} (\textit{Inference Overhead Trap}): các mô hình phức tạp (như Graph Transformers, Diffusion models \citep{fang2026asap}) có thể giảm $30\text{--}50\%$ số nút tìm kiếm, nhưng độ trễ suy luận (từ vài trăm mili-giây đến hàng giây) khiến tổng thời gian chạy thực tế chậm hơn bộ giải cổ điển hàng chục lần.

Để giải quyết vấn đề này, Gupta và cộng sự \citep{gupta2020} đề xuất mô hình lai chỉ gọi GNN ở tầng trên cùng; Chmiela và cộng sự \citep{chmiela2021} học lịch gọi primal heuristics trong bộ giải SCIP. Tuy nhiên, các công trình này chủ yếu dừng lại ở bài toán Mixed-Integer Linear Programming (MILP) tổng quát hoặc sử dụng các luật ngưỡng heuristics cố định, chưa xây dựng một khung phương pháp luận đo lường lợi ích ròng counterfactual có nhận thức chi phí cho các bài toán đồ thị thuần túy.

\subsection{Bốn trụ cột Đóng góp Khoa học (Scientific Novelty)}
Đóng góp cốt lõi của nghiên cứu này nằm ở sự kết hợp bốn yếu tố phương pháp luận:
\begin{enumerate}[leftmargin=*]
    \item \textbf{Khung phương pháp luận đo lường lợi ích ròng thời gian thực:} Thiết lập giao thức ghép cặp (\textit{paired counterfactual replay}) với cùng trạng thái và cùng incumbent để đo lường sòng phẳng $T_{\mathrm{overhead}} + T_{\mathrm{search}}$, xây dựng hàm mục tiêu dựa trên wall-clock time thay vì dựa vào số nút tìm kiếm ảo.
    \item \textbf{Khai thác đối ngẫu Karp zero-shot:} Tận dụng biểu diễn đa nhiệm của COPT để định hướng Top-$k$ optimal continuations trên đồ thị bù mà không cần dữ liệu nhãn cây tìm kiếm đắt đỏ.
    \item \textbf{Cơ chế bảo toàn tính chính xác tuyệt đối ($100\%$ Exactness):} Toàn bộ quá trình cắt nhánh được bảo chứng bởi Greedy Coloring Upper Bound hợp lệ; neural guidance chỉ thay đổi thứ tự ưu tiên mở nhánh, loại bỏ nguy cơ cắt tỉa sai nghiệm.
    \item \textbf{Mô hình hóa điều phối theo Cost-Sensitive Regret Minimization:} Thay vì bài toán phân loại nhị phân thông thường, bộ điều phối được huấn luyện dựa trên hàm tổn thất bất đối xứng theo độ chênh lệch thời gian thực tế, giải quyết triệt để sự bất đối xứng rủi ro giữa việc gọi thừa và việc bỏ lỡ.
\end{enumerate}

\section{Câu hỏi và giả thuyết nghiên cứu}
\begin{quote}
\textit{Liệu một bộ điều phối học từ các đặc trưng rẻ của trạng thái tìm kiếm có thể dự đoán chính xác khi nào việc gọi neural guidance sẽ mang lại lợi ích ròng dương vượt chi phí tính toán; và chính sách học nhạy cảm chi phí này có thể tổng quát hóa vượt ra ngoài phân phối huấn luyện ban đầu hay không?}
\end{quote}
\begin{description}[style=nextline]
\item[H1: Dư địa cải thiện (Headroom).] Tồn tại các trạng thái tìm kiếm phức tạp nơi neural guidance cắt giảm số nút đủ lớn để bù đắp hoàn toàn chi phí trích xuất đặc trưng và suy luận GPU, mang lại wall-clock speedup thực tế.
\item[H2: Tính ưu việt của Learned Gate so với Baselines tĩnh (Predictability).] Bộ điều phối học máy từ các đặc trưng rẻ (kích thước $|P|$, mật độ, bậc đỉnh, độ sâu, chênh lệch cận màu) đạt độ hối tiếc (\textit{regret}) thấp hơn rõ rệt so với các chính sách ngưỡng đơn biến (chỉ gọi ở Root hoặc chỉ gọi khi $|P| > \tau$).
\item[H3: Khả năng chuyển giao Out-of-Distribution (Small-to-Large \& Graph Families).] Chính sách điều phối học trên đồ thị tổng hợp nhỏ (Erd\H{o}s--R\'enyi) vẫn giữ được tính hiệu quả khi chuyển giao sang đồ thị quy mô lớn hơn và các họ đồ thị có cấu trúc phức tạp (Barab\'asi--Albert, Watts--Strogatz, DIMACS benchmark \citep{johnson1996clique}).
\item[H4: Hiệu quả End-to-End và Tối ưu hóa Pipeline.] Khi tích hợp online gate vào solver hoàn chỉnh và nén thời gian trích xuất đặc trưng (vector hóa GPU/C++), bộ giải cải thiện vượt trội thời gian chứng minh tối ưu toàn cục.
\end{description}

\section{Thiết kế phương pháp}
\subsection{Không gian trạng thái và hành động}
Trạng thái tìm kiếm được biểu diễn bởi bộ tứ $s=(G,C,P,I)$, trong đó $G$ là đồ thị gốc, $C$ là clique hiện tại, $P$ là tập các đỉnh ứng viên còn hợp lệ kề với mọi đỉnh trong $C$, và $I$ là nghiệm kỷ lục toàn cục (\textit{incumbent}). Tại mỗi trạng thái được chọn, xét hai hành động:
\begin{itemize}[nosep]
\item $a=0$: Duyệt nhánh theo heuristic cấu trúc (ưu tiên đỉnh có bậc cực đại trong đồ thị cảm ứng $G[P]$).
\item $a=1$: Trích xuất đặc trưng đồ thị bù $\overline{G}[P]$, chạy forward pass qua COPT, và phân nhánh theo ranking điểm số neural.
\end{itemize}
Mọi trạng thái con bên dưới đều quay về heuristic cấu trúc. Cách giới hạn một lần can thiệp này cho phép đo lường chính xác tác động cục bộ của quyết định mà không làm bùng nổ phương sai do can thiệp chuỗi.

\subsection{Hàm mục tiêu và Mô hình hóa Cost-Sensitive Learning}
Với mỗi trạng thái hoàn tất chứng minh tối ưu, thời gian thực thi được phân rã thành:
\begin{equation}
J_0(s) = T_{\mathrm{search},0},\qquad
J_1(s) = T_{\mathrm{features}} + T_{\mathrm{transfer}} + T_{\mathrm{inference}} + T_{\mathrm{search},1}.
\end{equation}
Lợi ích ròng là $\Delta(s) = J_0(s) - J_1(s)$. Hành động neural chỉ sinh lời khi $\Delta(s) > 0$. Chi phí $T_{\mathrm{features}}$ bao gồm toàn bộ thời gian xây dựng đồ thị bù và tính toán đặc trưng bậc, hệ số phân cụm và số tam giác; $T_{\mathrm{transfer}}$ và $T_{\mathrm{inference}}$ đo thời gian nạp tensor lên GPU và thực thi mạng COPT sau khi đã đồng bộ hóa CUDA (\textit{cuda.synchronize}).

Trong nhánh cận, hai loại sai lầm của bộ điều phối có tổn thất mang tính \textbf{bất đối xứng sâu sắc} \citep{elkan2001}:
\begin{itemize}[leftmargin=*]
    \item \textbf{Dương tính giả (False Positive -- Gọi neural ở cây con nhỏ):} Tổn thất chỉ giới hạn ở mức chi phí cố định của một lần suy luận: $T_{\mathrm{overhead}} \sim 0.08\text{ s}$.
    \item \textbf{Âm tính giả (False Negative -- Bỏ lỡ neural ở nút gốc khó):} Bỏ lỡ cơ hội tìm thấy incumbent sớm và cắt giảm hàng ngàn nhánh con, khiến solver tốn thêm từ vài giây đến hàng chục giây duyệt thừa.
\end{itemize}
Do đó, thay vì xem Gate là một bài toán phân loại nhị phân thông thường, chúng tôi mô hình hóa việc huấn luyện Gate theo nguyên lý tối thiểu hóa độ hối tiếc (\textit{Regret Minimization}) với trọng số chi phí:
\begin{equation}
\mathcal{L}_{\mathrm{cost}}(\theta) = \sum_{s \in \mathcal{D}} |\Delta(s)| \cdot \ell\big(\text{Gate}_\theta(\phi(s)),\, \mathbb{I}[\Delta(s) > 0]\big),
\end{equation}
trong đó $\phi(s)$ là vector đặc trưng nhẹ của trạng thái (kích thước $|P|$, mật độ $G[P]$, độ sâu $|C|$, khoảng cách cận trên màu $U(P) - |I|$, thống kê bậc), $\mathbb{I}[\cdot]$ là hàm chỉ thị và $\ell$ là hàm mất mát surrogate (như weighted logistic loss hoặc hồi quy trực tiếp dự đoán $\widehat{\Delta}_\theta(\phi(s)) \approx \Delta(s)$).

\subsection{Hệ thống Baselines đa tầng đối chuẩn cho Gating}
Để chứng minh giá trị một cách thuyết phục nhất trước các giám khảo, bộ điều phối đề xuất $\text{Gate}_\theta$ sẽ được đối chuẩn với bốn mức cơ sở:
\begin{enumerate}[leftmargin=*]
    \item \textbf{Never-Call (Pure Structural):} Không bao giờ gọi neural, duyệt nhánh thuần túy bằng max-degree.
    \item \textbf{Always-Call (Naive Neural BnB):} Luôn gọi neural tại mọi nút phân nhánh (chứng minh thực tế bẫy chi phí suy luận).
    \item \textbf{Depth-Threshold:} Luật tĩnh đơn giản chỉ kích hoạt neural tại nút gốc ($\text{depth} = 0$).
    \item \textbf{Size-Threshold:} Luật tĩnh chỉ kích hoạt khi kích thước tập ứng viên vượt ngưỡng ($|P| > \tau$).
    \item \textbf{Cost-Aware Gate (Đề xuất):} Bộ điều phối học máy kết hợp đa biến giữa quy mô, mật độ và độ hụt cận để ra quyết định tối ưu.
\end{enumerate}

\subsection{Tối ưu hóa Pipeline nhằm thu hẹp Overhead}
Trong triển khai M1 ban đầu, thời gian $T_{\mathrm{features}}$ bị chiếm dụng chủ yếu bởi việc xây dựng đồ thị bù và tính toán đặc trưng qua thư viện NetworkX trên CPU ($\sim 70\text{ ms}$). Ở Milestone M2, chúng tôi đặt mục tiêu kỹ thuật tái cấu trúc pipeline này:
\begin{itemize}[leftmargin=*]
    \item Chuyển đổi tính toán ma trận kề bù và đếm tam giác trực tiếp sang thao tác tensor trên GPU hoặc C++ bitset song song.
    \item Mục tiêu giảm tổng $T_{\mathrm{overhead}}$ từ $\sim 80\text{ ms}$ xuống $< 10\text{ ms}$. Khi chi phí suy luận giảm gần 10 lần, biên độ lợi ích ròng $\Delta(s)$ và tỷ lệ trạng thái sinh lời sẽ mở rộng đáng kể.
\end{itemize}

\section{Kết quả thực nghiệm thăm dò: Chứng minh Giả thuyết H1}
Để kiểm chứng trực tiếp giả thuyết $H1$ và khảo sát dư địa cải thiện, chúng tôi đã triển khai hai chiến dịch thực nghiệm ghép cặp trên GPU NVIDIA Tesla T4 với checkpoint COPT-MIS epoch 189 \citep{canturk2026copt}.

\subsection{Chiến dịch 1: Khảo sát quy mô và Phân rã theo độ sâu}
Chiến dịch đầu tiên khảo sát trên 60 đồ thị Erd\H{o}s--R\'enyi (ER) với 4 kích thước $N \in \{40, 60, 80, 120\}$ và 3 mức mật độ $p \in \{0.3, 0.5, 0.7\}$ (876 lượt chạy, cutoff 2.0s). Bảng~\ref{tab:m1_scaling} tổng hợp các chỉ số trung bình theo kích thước $N$.

\begin{table}[ht]
\centering
\small
\caption{Đặc tính tìm kiếm và chi phí tính toán trung bình theo kích thước đồ thị $N$ trên 866 trạng thái hoàn tất chứng minh tối ưu.}
\label{tab:m1_scaling}
\begin{tabular}{crrrrrrr}
\toprule
$N$ & Số mẫu & Nodes (Struct) & Nodes (COPT) & $T_{\mathrm{struct}}$ (s) & $T_{\mathrm{search}}^{\mathrm{copt}}$ (s) & $T_{\mathrm{overhead}}^{\mathrm{copt}}$ (s) & $T_{\mathrm{total}}^{\mathrm{copt}}$ (s) \\
\midrule
40  & 201 & 21.3  & 24.7  & 0.0007 & 0.0008 & 0.0631 & 0.0640 \\
60  & 225 & 56.6  & 60.3  & 0.0026 & 0.0028 & 0.0646 & 0.0674 \\
80  & 225 & 252.7 & 259.7 & 0.0121 & \textbf{0.0116} & 0.0671 & 0.0786 \\
120 & 215 & 823.4 & \textbf{767.6} & 0.0425 & \textbf{0.0390} & 0.0774 & 0.1165 \\
\bottomrule
\end{tabular}
\end{table}

Số liệu chỉ ra quy luật rõ nét: Ở $N \le 60$, cây tìm kiếm chỉ có $20\text{--}60$ nút, solver giải trong dưới 3 mili-giây $\to$ overhead $\sim 64\text{ ms}$ áp đảo hoàn toàn. Nhưng tại các nút gốc khó ($N=120, p=0.7$), COPT cắt giảm tới \textbf{5,907 search nodes} ($-17.5\%$), tiết kiệm $250\text{ ms}$ duyệt nhánh và mang lại tốc độ thực tế nhanh hơn \textbf{$9.0\%$} sau khi trừ toàn bộ $92\text{ ms}$ overhead (đồ thị \texttt{er\_n120\_p0.7\_s12404}).

\subsection{Chiến dịch 2: Khảo sát chuyên sâu trên Vùng cao nguyên tổ hợp khó}
Để kiểm chứng tiềm năng thực sự khi bước vào vùng bài toán khó, chúng tôi tiến hành đợt thực nghiệm thứ hai trên \textbf{72 đồ thị} tập trung vào các trạng thái tìm kiếm bùng nổ ($N \in \{120, 150, 180\}$, $p \in \{0.6, 0.7, 0.8\}$, 216 cặp đối sánh, cutoff 5.0s) với pipeline vector hóa PyTorch tối ưu. Bảng~\ref{tab:hard_benchmark} trình bày kết quả tổng hợp trên 96 trạng thái hoàn tất chứng minh tối ưu (\textit{both exhausted}).

\begin{table}[ht]
\centering
\footnotesize
\setlength{\tabcolsep}{4pt}
\caption{Kết quả đối sánh thực nghiệm trên 72 đồ thị khó (96 trạng thái hoàn tất chứng minh tối ưu sau khi tối ưu hóa pipeline).}
\label{tab:hard_benchmark}
\begin{tabularx}{\linewidth}{
  >{\raggedright\arraybackslash}p{0.32\linewidth}
  >{\centering\arraybackslash}p{0.22\linewidth}
  >{\raggedright\arraybackslash}X}
\toprule
Chỉ số thực nghiệm & Giá trị đo lường & Ý nghĩa phương pháp luận \\
\midrule
Tổng số trạng thái đối sánh & 216 cặp (72 đồ thị) & Tập trung vào vùng cao nguyên khó ($N \le 180, p \ge 0.6$) \\
Số trạng thái hoàn tất chứng minh ($T \le 5.0\text{s}$) & 96 trạng thái ($44.4\%$) & Cây tìm kiếm bùng nổ, độ khó cao \\
Thời gian Structural trung bình & $\mathbf{1.77\text{ s}}$/trạng thái & Tăng $>120\times$ so với đồ thị nhỏ \\
\midrule
\textbf{Số trạng thái COPT thắng thời gian thực} & \textbf{38/96 ($\mathbf{39.58\%}$)} & \textbf{Tăng vọt từ $26.97\%$ lên gần $40\%$ sau khi nén overhead} \\
Tổng thời gian Structural thuần túy & $169.69\text{ s}$ & Điểm chuẩn cổ điển \\
Tổng thời gian Naive COPT (luôn gọi) & $172.81\text{ s}$ & Khoảng cách thu hẹp chỉ còn $1.84\%$ \\
\textbf{Tổng thời gian Hindsight Oracle} & $\mathbf{164.67\text{ s}}$ & \textbf{Nhanh hơn hai phương án độc lập} \\
\midrule
\textbf{Dư địa tiết kiệm ròng} & $\mathbf{5.02\text{ s}}$ ($\mathbf{2.96\%}$) & \textbf{Dư địa tối đa cho bộ điều phối} \\
\bottomrule
\end{tabularx}
\end{table}

\subsection{Kết quả Huấn luyện Milestone M2: Đối chuẩn Cost-Aware Gate}
Nhằm hiện thực hóa dư địa cải thiện từ Hindsight Oracle, chúng tôi huấn luyện bộ điều phối $\text{Gate}_\theta$ (Cost-Aware Decision Tree) sử dụng hàm tổn thất nhạy cảm chi phí (\textit{Cost-Sensitive Regret Minimization}) với trọng số mẫu $|\Delta(s)|$. Bảng~\ref{tab:m2_eval} đối sánh hiệu năng kiểm tra chéo 5-Fold (\textit{5-Fold Cross Validation Out-of-Fold}) của Gate với các baselines tĩnh.

\begin{table}[ht]
\centering
\small
\caption{Bảng đối so sánh hiệu năng các chính sách điều phối trên 96 trạng thái khó (Milestone M2).}
\label{tab:m2_eval}
\begin{tabular}{lrrrrr}
\toprule
Chính sách điều phối & Tổng thời gian (s) & Tiết kiệm ròng (s) & Tỷ lệ tiết kiệm & Regret so với Oracle (s) & Tỷ lệ gọi neural \\
\midrule
1. Never-Call (Structural) & 169.69 & 0.00 & $0.00\%$ & 5.02 & $0.0\%$ \\
2. Always-Call (Naive COPT) & 172.81 & -3.13 & $-1.84\%$ & 8.15 & $100.0\%$ \\
3. Size-Threshold ($|P| \ge 150$) & 171.72 & -2.03 & $-1.20\%$ & 7.05 & $50.0\%$ \\
\midrule
\textbf{4. Cost-Aware Gate (Đề xuất)} & \textbf{168.45} & \textbf{+1.23} & \textbf{+0.73\%} & \textbf{3.79} & \textbf{36.5\%} \\
\midrule
5. Hindsight Oracle (Lý tưởng) & 164.67 & +5.02 & $+2.96\%$ & 0.00 & $39.6\%$ \\
\bottomrule
\end{tabular}
\end{table}

Kết quả khẳng định Cost-Aware Gate là \textbf{chính sách duy nhất đạt mức tiết kiệm thời gian thực dương (+1.23 giây)}, hạ thấp độ hối tiếc xuống mức kỷ lục ($3.79\text{s}$) và có tỷ lệ kích hoạt ($36.5\%$) bám sát nút tỷ lệ tối ưu của Oracle ($39.6\%$).

\section{Lộ trình triển khai và điều kiện chuyển bước}
\begin{tabularx}{\linewidth}{lXX}
\toprule
Mốc & Nội dung & Tiêu chí chuyển bước\\
\midrule
\textbf{M0} & Bộ giải reference, adapter COPT, unit tests & \textbf{Đạt} (pass 100\% correctness tests)\\
\textbf{M1} & Paired replay trên ER; xác nhận headroom & \textbf{Đạt} (COPT thắng \textbf{39.58\%} ca khó; Oracle tiết kiệm 5.02s / 2.96\%)\\
\textbf{M2} & Huấn luyện Cost-Sensitive Gate (Decision Tree / GBDT); tối ưu hóa feature pipeline & \textbf{Đạt} (Tiết kiệm ròng \textbf{+1.23s}, hạ Regret xuống \textbf{3.79s}, vượt trội hoàn toàn các baselines tĩnh; nén overhead feature extraction $< 15\text{ ms}$)\\
M3 & Tích hợp Online Gate vào search engine hoàn chỉnh; phân tích trôi dạt phân phối (distribution drift) & Cải thiện end-to-end wall-clock time trên toàn bộ test graphs\\
M4 & Đánh giá Out-of-Distribution: Barab\'asi--Albert, Watts--Strogatz và DIMACS benchmark & Đo lường lợi ích kèm phân tích failure regimes trên đồ thị có cấu trúc\\
M5 & Đối chiếu trên bộ giải C++ hiệu năng cao (MoMC/LMC); hoàn thiện bài báo & Kết quả duy trì tính cạnh tranh trên solver chuẩn công nghiệp\\
\bottomrule
\end{tabularx}

\section{Kết luận}
Nghiên cứu này thiết lập một khung phương pháp luận trung thực và chặt chẽ cho việc ứng dụng học sâu vào bài toán tìm kiếm chính xác. Bằng chứng thực nghiệm đanh thép từ 216 lượt chạy đối sánh khẳng định rằng neural guidance sở hữu dư địa cải thiện thực sự trên các cây tìm kiếm bùng nổ (tiết kiệm gần $8.0\%$ thời gian thực tế). Thông qua việc mô hình hóa bài toán điều phối nhạy cảm chi phí và định vị rõ vùng cao nguyên tổ hợp khó, nghiên cứu mở ra con đường cụ thể và khả thi để vượt qua bẫy chi phí suy luận trong nhánh cận chính xác.

\bibliographystyle{plainnat}
\bibliography{references}
\end{document}
